This section displays the scatter plots and cluster-example grids generated by
\texttt{main.py}, mirroring the figure gallery in the Jupyter notebook.

\subsection{Scatter Plots}

\begin{figure}[H]
  \centering
  \includegraphics[width=0.75\linewidth]{scatter_Raw.png}
  \caption{Cluster assignments on raw 784-D pixels (2-D reference view). Clusters overlap
    massively; the algorithm groups points by overall brightness and silhouette area rather than
    garment type.}
\end{figure}

\begin{figure}[H]
  \centering
  \includegraphics[width=0.75\linewidth]{scatter_PCA.png}
  \caption{Cluster assignments on 85-D PCA features. Slightly tidier than Raw — PCA strips
    isotropic pixel noise — but the class manifolds remain entangled.}
\end{figure}

\begin{figure}[H]
  \centering
  \includegraphics[width=0.75\linewidth]{scatter_SAE.png}
  \caption{Cluster assignments on the 32-D dense-SAE bottleneck. Some local structure emerges,
    but several clusters still span multiple true labels.}
\end{figure}

\begin{figure}[H]
  \centering
  \includegraphics[width=0.75\linewidth]{scatter_CNN-SAE.png}
  \caption{Cluster assignments on CNN-SAE features. Boundaries between clusters are sharper,
    especially for footwear, matching the best Adjusted Rand among autoencoder methods.}
\end{figure}

\begin{figure}[H]
  \centering
  \includegraphics[width=0.75\linewidth]{scatter_t-SNE.png}
  \caption{Cluster assignments on the 2-D t-SNE embedding. Clusters pop out as well-defined
    blobs, explaining the order-of-magnitude jump in Silhouette and Calinski--Harabasz. Global
    distances are not meaningful in t-SNE.}
\end{figure}

\begin{figure}[H]
  \centering
  \includegraphics[width=0.75\linewidth]{scatter_UMAP.png}
  \caption{Cluster assignments on the 2-D UMAP embedding — our best overall configuration.
    Clusters are compact, well-separated, and global topology is better preserved than in t-SNE.}
\end{figure}

\subsection{Cluster Examples}

\begin{figure}[H]
  \centering
  \includegraphics[width=0.85\linewidth]{cluster_examples_Raw.png}
  \caption{Representative images from each cluster on Raw features. Clusters are not
    class-aligned — a single cluster often mixes T-shirts, shirts, and pullovers.}
\end{figure}

\begin{figure}[H]
  \centering
  \includegraphics[width=0.85\linewidth]{cluster_examples_PCA.png}
  \caption{Cluster examples after PCA. Footwear clusters begin to cohere, but upper-body
    garments are still routinely co-clustered.}
\end{figure}

\begin{figure}[H]
  \centering
  \includegraphics[width=0.85\linewidth]{cluster_examples_SAE.png}
  \caption{Cluster examples after the fully-connected SAE. Shoe and trouser clusters are now
    clearly homogeneous; bag clusters become distinct. Shirts / pullovers / coats still bleed
    into each other.}
\end{figure}

\begin{figure}[H]
  \centering
  \includegraphics[width=0.85\linewidth]{cluster_examples_CNN-SAE.png}
  \caption{Cluster examples after CNN-SAE. Per-cluster images are visibly more consistent;
    footwear subtypes begin to split (sneakers vs.\ ankle-boots), and dresses / trousers form
    clean clusters.}
\end{figure}

\begin{figure}[H]
  \centering
  \includegraphics[width=0.85\linewidth]{cluster_examples_t-SNE.png}
  \caption{Cluster examples after t-SNE. Each cluster reads as essentially one garment type.
    Residual ambiguity is confined to shirt $\leftrightarrow$ pullover $\leftrightarrow$ coat.}
\end{figure}

\begin{figure}[H]
  \centering
  \includegraphics[width=0.85\linewidth]{cluster_examples_UMAP.png}
  \caption{Cluster examples after UMAP — the cleanest result of the study. The 10 clusters
    correspond almost one-to-one to the 10 Fashion-MNIST classes, with residual confusion
    confined to shirt $\leftrightarrow$ pullover $\leftrightarrow$ coat.}
\end{figure}
